\newlabel{fig:1}{{1}{4}{\textbf {Overview of the DIReCT$++$ framework for multimodal PET synthesis and downstream analysis.} (A) Tracer-conditioned rectified flow from MRI to FDG- or AV45-PET. (B) Lightweight affine fine-tuning of tracer guidance derived from frozen BiomedCLIP encoders. The symbols $\alpha $ and $\beta $ denote FDG adapter parameters $(\alpha _1,\alpha _2)$ and AV45 adapter parameters $(\beta _1,\beta _2)$, respectively; these adaptations are separate from the raw tracer embeddings used by the reconstructed Full generation experiments. (C) A shared conditional 3D U-Net predicts velocity fields for the requested tracer. (D) Learned one-step PET estimation by residual addition of the predicted velocity to MRI; each forward pass produces one tracer. (E) Image-quality, regional uptake and cross-sectional classification analyses.\relax }{figure.caption.1}{}}
\newlabel{tab:historical_generation}{{1}{5}{Absolute historical Figure 2 summary values\relax }{table.1}{}}
\newlabel{fig:2}{{2}{7}{\textbf {Historical PET synthesis comparisons and illustrative examples.} (A) Radar plots summarize the historical aggregate SSIM, MSE, MAE and PSNR values for ADNI FDG/AV45 and OASIS-3 AV45. Metric values are displayed relative to the best value among methods; absolute aggregate values and scaling conventions are listed in Table~\ref {tab:historical_generation}. These aggregates are separate from the reconstructed OOF and masked analyses. (B) Illustrative ADNI CN, MCI and AD examples with MRI, competing synthesis outputs and real PET; paired subrows show FDG and AV45. (C) Historical CN/AD OASIS-3 examples are limited to AV45. Examples were randomly selected within diagnostic groups as reported by the authors; the selection seed could not be verified from the available records. This differs from the supplemental quantile-based selection and does not establish population-level superiority.\relax }{figure.caption.2}{}}
\newlabel{fig:3}{{3}{8}{\textbf {Detailed regional visualization of synthetic and real PET images.} The figure displays magnified regional views of FDG-PET (left two columns) and AV45-PET (right two columns) for representative AD and CN subjects. For each modality and group, the first sub-column shows real PET, and the second sub-column shows synthetic PET. Key anatomical regions are highlighted in the bottom-right corner of each sub-column, with their names labeled at the top-center. Each row corresponds to a distinct brain region. Displayed uptake values are arbitrary units after per-image min--max normalization, without anatomical reference-region normalization. Cases were randomly selected within their diagnostic groups as reported by the authors; the selection seed could not be verified from the available records. These selected cases illustrate regional appearance and do not establish population-level agreement or clinical SUV/SUVR calibration. \relax }{figure.caption.3}{}}
\newlabel{fig:4}{{4}{10}{\textbf {Regional comparison of mean normalized uptake between real and MRI-synthetic PET images.} (A,B) Violin plots show regional mean normalized uptake (arbitrary units) across cortical and subcortical regions for real and synthetic PET, comparing AV45 (A) and FDG (B). All regional uptake values were derived from per-image min--max-normalized PET intensities without anatomical reference-region normalization. Panels A,B summarize the pooled available subjects without diagnostic stratification. Non-significant differences are not evidence of equivalence. (C--F) AD--CN regional comparisons. Violin plots display mean normalized uptake in AD (red) and CN (green) subjects for selected regions. Each region uses paired half-violins: upper (solid) = real PET, lower (hatched) = synthetic PET; colors match across modalities. Solid lines indicate real PET significance, dashed lines indicate synthetic PET significance. The statistical annotations in panels (C--F) follow the BH-FDR analysis: \texttt {****} corresponds to $q<10^{-4}$ and \texttt {ns} to $q\geq 0.05$. (C,D) show representative regions with significant AD--CN differences for AV45 and FDG, whereas (E,F) show selected regions without significant AD--CN differences after FDR correction. Cblm-WM denotes Cerebellum-White-Matter. Additional regional agreement and effect-size analyses are provided in the Supplementary Materials. \relax }{figure.caption.4}{}}
\newlabel{fig:5}{{5}{12}{\textbf {Cross-sectional classification under the revised three-fold protocol.} Rows A--C correspond to AD/CN, EMCI+LMCI/CN and EMCI/LMCI. Within each row, the five metric columns compare nine input conditions: MRI; MRI+Real FDG; MRI+Syn FDG; MRI+Real AV45; MRI+Syn AV45; MRI+Real-both; MRI+Syn-both; Clinical-only; and MRI+Clinical. Points are equally weighted three-fold means; horizontal error bars show sample s.d. in percentage points, not 95\% CIs. ACC, AUC, SEN, SPE and balanced ACC are reported with their CIs in Table~\ref {tab:classification_summary}. The results do not demonstrate longitudinal prognosis, PET equivalence or benefit independent of diagnosis-related clinical inputs.\relax }{figure.caption.5}{}}
\newlabel{eq:align_opt}{{1}{15}{Domain-Adapted VLM}{equation.4.1}{}}
\newlabel{eq:rf_ode}{{2}{16}{Rectified Flow}{equation.4.2}{}}
\newlabel{eq:composite_conditioning}{{3}{16}{Rectified Flow}{equation.4.3}{}}
\newlabel{eq:rf_regression}{{4}{16}{Rectified Flow}{equation.4.4}{}}
\newlabel{eq:rf_loss}{{5}{16}{Rectified Flow}{equation.4.5}{}}
\newlabel{eq:multitask_loss}{{6}{17}{Multi-Task Learning and One-Step Generation}{equation.4.6}{}}
\newlabel{eq:onestep_generation}{{7}{17}{Multi-Task Learning and One-Step Generation}{equation.4.7}{}}
